% Generated draft: requires mathematical and source review.
\documentclass[11pt]{article}
\usepackage[margin=25mm]{geometry}
\usepackage{fontspec}
\setmainfont{Latin Modern Roman}
\usepackage{amsmath,amssymb,mathtools,unicode-math}
\setmathfont{Latin Modern Math}
\usepackage{enumitem,needspace,xurl,hyperref}
\setlength{\parindent}{0pt}
\setlength{\parskip}{5pt}
\setlength{\emergencystretch}{3em}
\newcommand{\N}{\mathbb{N}}
\newcommand{\Z}{\mathbb{Z}}
\newcommand{\Q}{\mathbb{Q}}
\newcommand{\R}{\mathbb{R}}
\newcommand{\C}{\mathbb{C}}
\newcommand{\F}{\mathbb{F}}
\newcommand{\A}{\mathbb{A}}
\newcommand{\PP}{\mathbb{P}}
\newcommand{\E}{\mathbb{E}}
\newcommand{\eps}{\varepsilon}
\newcommand{\dd}{\,\mathrm{d}}
\begin{document}
\section*{UnsolvedMath ID: 2307010}
\textbf{Source problem code:} catalogue entry
\subsection{Definitions and mathematical statement}
Let $\Omega,\Omega'\subset\C$ be plane domains and $K>1$. A $K$-quasiconformal map is an orientation-preserving homeomorphism $f:\Omega\to\Omega'$ in $W^{1,2}_{\mathrm{loc}}$ satisfying
\[
 |f_{\bar z}|\leq\frac{K-1}{K+1}|f_z|\quad\text{almost everywhere}.
\]
The requested assertion is that all real partial derivatives of $f$ belong to $L^p_{\mathrm{loc}}(\Omega)$ for every
\[
 2\leq p<2+\frac2{K-1}=\frac{2K}{K-1}.
\]
Equivalently, for every compact $E\subset\Omega$ and every such $p$, $\int_E|Df|^p<\infty$. No endpoint inclusion is asserted. The source's phrase locally $L$-integrable is read in its displayed $p$-range as local $L^p$ integrability. For $K=1$ the written coefficient is undefined; the usual conformal limiting case is distinct from the $K>1$ formula and is not silently assigned a finite threshold.

The source attributes a weaker positive improvement $2\leq p<2+c(K)$ to Boyarskiĭ and asks for $c(K)=2/(K-1)$. It gives the radial example
\[
 f(z)=|z|^{1/K-1}z\quad(z\ne0),\qquad f(0)=0,
\]
to show sharpness: its derivative has the radial scale $|z|^{1/K-1}$, so the critical exponent is precisely $2K/(K-1)$.

The supplied assessment attributes resolution of the planar sharp range to Astala. Its discussion of unresolved sharp constants in dimensions $n\geq3$ concerns an extension absent from this source question and is not imported into its quantifiers. The summary's reference to $K/(K-1)$-related exponents must be distinguished from the displayed derivative threshold, which is $2K/(K-1)$.
\subsection{Short English statement}
Prove the sharp higher-integrability range for derivatives of a planar quasiconformal map. The supplied registry labels the question solved and its summary attributes the planar result to Astala, while describing higher-dimensional extensions as partial/open. Those extensions are not part of the literal planar problem. The radial example fixes the excluded endpoint. No new proof or independent literature verification is claimed.
\subsection{Sources}
\begin{itemize}
\item UnsolvedMath permanent ID 2307010, source code catalogue entry, ``Research Problems in Function Theory --- Problem 7.10.'' Registry curation and metadata: Ulam AI and UnsolvedMath Contributors, CC BY 4.0; updated 2026-08-06.
\item Hayman, \emph{Research Problems in Function Theory} (2018 source update), Problem 7.10, attributed to F. W. Gehring: \url{https://arxiv.org/abs/1809.07200}. Supplied source-TeX extraction accessed 2026-07-29. The metadata retains the update for status triage and marks status/release and rights review NEEDS\_REVIEW.
\item Boyarskiĭ is named for the initial integrability improvement, and Astala for the reported sharp planar theorem. Complete titles, dates and theorem locators are not supplied; none are inferred.
\end{itemize}
\end{document}
